% The content of the blueprint, shared by the web and the print version.
% One chapter per theme, one section per theorem (so the web version has a
% page per theorem).

\chapter*{How to read this}

Atlas couples solvers, classical or learned, through typed ports, and its
proposal makes claims about convergence and error.  This blueprint lists the
mathematical statements those claims rest on.  Each statement appears three
times:

\begin{itemize}
  \item \textbf{in plain words}, so a reader can judge whether it says what
    the proposal needs;
  \item \textbf{as mathematics}, with a proof on paper;
  \item \textbf{as a Lean 4 declaration}, checked by the Lean kernel against
    Mathlib.  A statement counts as machine-checked only when it has no
    \texttt{sorry} and \texttt{\#print axioms} shows nothing beyond Lean's
    three standard axioms (\texttt{propext}, \texttt{Classical.choice},
    \texttt{Quot.sound}).
\end{itemize}

A machine check guarantees that the proof proves the statement.  It does not
guarantee that the statement is the one that was meant, which is why the plain
words are here.

Everything in this blueprint is either a statement about metric and normed
spaces that holds in any dimension, or finite-dimensional linear algebra.  The
continuous theory of partial differential equations (Sobolev spaces, traces) is
cited by the proposal, not formalised here.

The labels T1, T3, \dots{} are those of the project's formal-proofs plan.

\chapter{Contractions: convergence and certificates}

A coupled iteration is a map $T$ on the data exchanged at the interfaces: one
sweep sends the current interface data $a$ to $T(a)$, and the coupled answer is
a point the sweep leaves unchanged, $T(a^\star)=a^\star$.  The two statements of
this chapter are the two questions one can ask of such a map.  How far is the
answer of a learned iteration from the classical one (T1)?  And how far is the
state I am holding from the answer (T7)?  They are the same inequality read two
ways.

\section{T1: a perturbed contraction}

\paragraph{In plain words.}
Replace the classical coupling map by a learned one.  If the learned map shrinks
every distance by a factor $\rho<1$, and at the classical answer it differs from
the classical map by at most $\delta$, then the learned iteration has exactly
one answer, it reaches that answer geometrically from any starting guess, and
that answer is within $\delta/(1-\rho)$ of the classical one.  Nothing is
assumed about how the learned map was trained.  The classical map need not be a
contraction; it only needs to have an answer.

\begin{theorem}[T1, perturbed contraction]\label{thm:T1}
  \lean{Atlas.perturbed_contraction}
  \leanok
  Let $(X,d)$ be a complete metric space and $T_c,T_l:X\to X$.  Suppose
  \begin{enumerate}
    \item $T_c(a_c)=a_c$ for some $a_c\in X$;
    \item $d(T_l x,T_l y)\le\rho\,d(x,y)$ for all $x,y\in X$, with $0\le\rho<1$;
    \item $d(T_l a_c,\,T_c a_c)\le\delta$.
  \end{enumerate}
  Then there is a point $a_l\in X$ such that
  \begin{enumerate}
    \item $T_l(a_l)=a_l$, and $a_l$ is the only fixed point of $T_l$;
    \item $d(T_l^{\,k}x_0,\,a_l)\le\rho^{k}\,d(x_0,a_l)$ for every $x_0\in X$ and every $k\in\N$;
    \item $d(a_l,a_c)\le\dfrac{\delta}{1-\rho}$.
  \end{enumerate}
\end{theorem}

\begin{proof}
  Existence and uniqueness of $a_l$ are Banach's fixed-point theorem
  (Mathlib: \texttt{ContractingWith.fixedPoint}).  For the rate,
  $T_l^{\,k}a_l=a_l$, so
  $d(T_l^{\,k}x_0,a_l)=d(T_l^{\,k}x_0,T_l^{\,k}a_l)\le\rho^k d(x_0,a_l)$ by
  applying hypothesis 2 $k$ times.  For the distance, the triangle inequality
  gives
  \[
    d(a_l,a_c)\le d(T_la_l,T_la_c)+d(T_la_c,a_c)\le\rho\,d(a_l,a_c)+\delta,
  \]
  using $T_ca_c=a_c$ in hypothesis 3, and $1-\rho>0$.
\end{proof}

\begin{corollary}[T1 with the defect bounded everywhere]\label{cor:T1-uniform}
  \lean{Atlas.perturbed_contraction_uniform}
  \leanok
  \uses{thm:T1}
  The conclusion of Theorem~\ref{thm:T1} holds if hypothesis 3 is replaced by
  $d(T_lx,T_cx)\le\delta$ for all $x\in X$.
\end{corollary}

\begin{proof}
  \uses{thm:T1}
  Take $x=a_c$.
\end{proof}

\begin{proposition}[T1 is sharp]\label{prop:T1-sharp}
  \lean{Atlas.perturbed_contraction_sharp}
  \leanok
  For every $\rho\in[0,1)$ and $\delta\ge0$ there are maps $T_c,T_l:\R\to\R$
  satisfying the hypotheses of Corollary~\ref{cor:T1-uniform} whose fixed points
  are exactly $\delta/(1-\rho)$ apart.
\end{proposition}

\begin{proof}
  $T_c(x)=\rho x$ and $T_l(x)=\rho x+\delta$, with fixed points $0$ and
  $\delta/(1-\rho)$.
\end{proof}

\section{T7: the a-posteriori certificate}

\paragraph{In plain words.}
A solver hands back a state $w$.  Its residual, how far one more classical step
moves it, can be computed.  Its error, how far it is from the settled state
$w^\star$, cannot, because $w^\star$ is unknown.  The certificate bounds the
second by the first: if the step brings $w$ closer to $w^\star$ by a factor
$L<1$, the error is at most the residual divided by $1-L$.

\paragraph{What the constant must be.}
The factor $L$ must hold \emph{at the state being certified}.  A constant read
off one approach to $w^\star$ says nothing about a state that approaches from
another direction.  For a linear step this is exact: the smallest constant that
certifies every state is an operator norm, a supremum over all directions
(Theorem~\ref{thm:T7-opnorm}).  The project met this the hard way: a constant
read off one march under-bounded two of seven solver arms on its largest test.

\begin{theorem}[T7, the certificate]\label{thm:T7}
  \lean{Atlas.certificate}
  \leanok
  Let $(X,d)$ be a metric space, $\Phi:X\to X$, and $w,w^\star\in X$ with
  $\Phi(w^\star)=w^\star$.  If $d(\Phi w,\Phi w^\star)\le L\,d(w,w^\star)$ with
  $L<1$, then
  \[
    d(w,w^\star)\le\frac{d(w,\Phi w)}{1-L}.
  \]
\end{theorem}

\begin{proof}
  $d(w,w^\star)\le d(w,\Phi w)+d(\Phi w,\Phi w^\star)\le d(w,\Phi w)+L\,d(w,w^\star)$,
  and $1-L>0$.
\end{proof}

\begin{corollary}[T7, Banach's form]\label{cor:T7-banach}
  \lean{Atlas.certificate_of_contraction}
  \leanok
  \uses{thm:T7}
  Let $X$ be a non-empty complete metric space and
  $d(\Phi x,\Phi y)\le L\,d(x,y)$ for all $x,y$, with $0\le L<1$.  Then $\Phi$
  has a unique fixed point $w^\star$, and
  $d(w,w^\star)\le d(w,\Phi w)/(1-L)$ for every $w\in X$.
\end{corollary}

\begin{proof}
  \uses{thm:T7}
  Banach's theorem gives $w^\star$; Theorem~\ref{thm:T7} applies to every $w$.
  (Mathlib: \texttt{ContractingWith.dist\_fixedPoint\_le}.)
\end{proof}

\begin{theorem}[T7 for a linear step: the constant is an operator norm]\label{thm:T7-opnorm}
  \lean{Atlas.certificate_constant_is_opNorm}
  \leanok
  Let $E$ be a real normed space, $A:E\to E$ bounded and linear, $b\in E$, and
  $\Phi(w)=Aw+b$.  Suppose $I-A$ has a bounded inverse $B$, and
  $\Phi(w^\star)=w^\star$.  Then
  \begin{enumerate}
    \item $\norm{w-w^\star}\le\norm{B}\,\norm{\Phi(w)-w}$ for every $w\in E$;
    \item if $C\ge0$ and $\norm{w-w^\star}\le C\,\norm{\Phi(w)-w}$ for every
      $w\in E$, then $\norm{B}\le C$.
  \end{enumerate}
  So $\norm{(I-A)^{-1}}$ is the least constant that certifies every state.
\end{theorem}

\begin{proof}
  $\Phi(w)-w=(A-I)(w-w^\star)$, so $w-w^\star=-B(\Phi(w)-w)$, which gives 1.
  For 2, given $x\in E$ let $w=w^\star+Bx$.  Then
  $\Phi(w)-w=(A-I)Bx=-x$, and the hypothesis reads $\norm{Bx}\le C\norm{x}$.
  This holds for every $x$, so $\norm{B}\le C$.
\end{proof}

\begin{example}[The two-mode example]\label{ex:T7-two-mode}
  \lean{Atlas.certificate_two_mode_example}
  \leanok
  For $\Phi(x,y)=(0.6\,x,\ 0.95\,y)$ on $\R^2$ with the maximum norm, whose
  settled state is $0$: the error of $(1,0)$ is $2.5$ times its residual, and
  the error of $(0,1)$ is $20$ times its residual.  A constant of $2.5$, read
  off a march that stays in the fast mode, under-bounds the slow mode by a
  factor of eight.
\end{example}

\begin{remark}
  T1's distance bound is T7 applied to the learned map at the classical answer:
  the residual of $a_c$ under $T_l$ is at most $\delta$.
\end{remark}

\chapter{The master error bound}

A simulation marches $u^{n+1}=\Phi(u^n)$, where $\Phi$ is one whole coupled
macro-step.  The true solution sampled at the same instants is $u^{n,\star}$.
This chapter bounds the error $e^n=u^n-u^{n,\star}$ after $N$ steps.

The law is the same for a single monolithic network and for a coupled system.
What a decomposition changes is the constants: the one-step defect gains two
terms from the coupling, and the stability constant $L$ factors over the parts.

\section{T4: recursion, split, accumulation, regimes}

\paragraph{In plain words.}
\begin{itemize}
  \item \emph{Recursion.}  The new error is the old error carried through one
    step, plus the one-step defect: what a single step does wrong when it starts
    from the true solution.
  \item \emph{Three-term split.}  With a coupling, the defect is a sum of three
    parts: the experts being wrong even when handed the true interface data
    ($\tau$), the interface problem being posed with the wrong operator
    ($\sigma$), and the interface solve being stopped early ($\gamma$).
    Iterating the coupling reduces $\gamma$ and nothing else.
  \item \emph{Accumulation.}  If one step magnifies a difference by at most $L$,
    the error after $N$ steps is at most $L^N$ times the initial error, plus
    every defect magnified by the steps that follow it.
  \item \emph{Three regimes.}  With every defect at most $\delta$: the error
    stays below $\delta/(1-L)$ for all time when $L<1$; it grows at most
    linearly, $N\delta$, when $L=1$; the bound grows exponentially when $L>1$.
\end{itemize}

\begin{lemma}[T4a, the exact error recursion]\label{lem:T4-recursion}
  \lean{Atlas.error_recursion}
  \leanok
  Let $E$ be a normed space, $\Phi:E\to E$, and $u,u^\star:\N\to E$ with
  $u^{n+1}=\Phi(u^n)$.  Then for every $n$
  \[
    u^{n+1}-u^{n+1,\star}
      =\bigl[\Phi(u^n)-\Phi(u^{n,\star})\bigr]
      +\bigl[\Phi(u^{n,\star})-u^{n+1,\star}\bigr].
  \]
  The second bracket is the one-step defect $d^{n+1}$.
\end{lemma}

\begin{proof}
  Add and subtract $\Phi(u^{n,\star})$.
\end{proof}

\begin{lemma}[T4b, the three-term split]\label{lem:T4-split}
  \lean{Atlas.defect_split}
  \leanok
  Let $\Phi_\mu:E\to E$ be the composed step forced to use the interface datum
  $\mu$, and $\mathcal S:E\to E$ the exact evolution.  For any three data
  $\lambda^\star$, $\lambda^\dagger$, $\lambda^{(k)}$ and any state $v$,
  \[
    \Phi_{\lambda^{(k)}}(v)-\mathcal S(v)
    =\underbrace{\Phi_{\lambda^\star}(v)-\mathcal S(v)}_{\tau}
    +\underbrace{\Phi_{\lambda^\dagger}(v)-\Phi_{\lambda^\star}(v)}_{\sigma}
    +\underbrace{\Phi_{\lambda^{(k)}}(v)-\Phi_{\lambda^\dagger}(v)}_{\gamma}.
  \]
\end{lemma}

\begin{proof}
  The sum telescopes.
\end{proof}

\begin{lemma}[T4c, accumulation]\label{lem:T4-accumulation}
  \lean{Atlas.accumulation}
  \leanok
  Let $L\ge0$ and $e,d:\N\to\R$ with $e_{n+1}\le L\,e_n+d_{n+1}$ for all $n$.
  Then for every $N$
  \[
    e_N\le L^N e_0+\sum_{n=1}^{N}L^{N-n}d_n .
  \]
\end{lemma}

\begin{proof}
  Induction on $N$.  The step multiplies the bound for $N$ by $L\ge0$ and adds
  $d_{N+1}$.
\end{proof}

\begin{theorem}[T4, the master bound]\label{thm:T4}
  \lean{Atlas.master_bound}
  \leanok
  \uses{lem:T4-recursion, lem:T4-accumulation}
  Let $u^{n+1}=\Phi(u^n)$, and suppose
  $\norm{\Phi(u^n)-\Phi(u^{n,\star})}\le L\,\norm{u^n-u^{n,\star}}$ for every
  $n$, with $L\ge0$.  Then for every $N$
  \[
    \norm{u^N-u^{N,\star}}\le L^N\norm{u^0-u^{0,\star}}
      +\sum_{n=1}^{N}L^{N-n}\,\norm{\Phi(u^{n-1,\star})-u^{n,\star}} .
  \]
  The stability hypothesis is needed only along the two trajectories.
\end{theorem}

\begin{proof}
  \uses{lem:T4-recursion, lem:T4-accumulation}
  Take norms in Lemma~\ref{lem:T4-recursion} and apply
  Lemma~\ref{lem:T4-accumulation} to $e_n=\norm{u^n-u^{n,\star}}$.
\end{proof}

\begin{corollary}[T4 with the three terms named]\label{cor:T4-three}
  \lean{Atlas.master_bound_three_terms}
  \leanok
  \uses{thm:T4, lem:T4-split}
  Let the composed step be $\Phi(v)=\Phi_{\lambda^{(k)}(v)}(v)$, where
  $\lambda^{(k)}(v)$ is the datum the interface solver returns at the state $v$.
  Under the hypotheses of Theorem~\ref{thm:T4}, for any sequences of data
  $\lambda^{\star}_n$, $\lambda^{\dagger}_n$,
  \[
    \norm{u^N-u^{N,\star}}\le L^N\norm{u^0-u^{0,\star}}
      +\sum_{n=1}^{N}L^{N-n}\bigl(\tau^n+\sigma^n+\gamma^n\bigr),
  \]
  with $\tau^n$, $\sigma^n$, $\gamma^n$ the norms of the three terms of
  Lemma~\ref{lem:T4-split} at $v=u^{n-1,\star}$ and
  $\mathcal S(u^{n-1,\star})=u^{n,\star}$.
\end{corollary}

\begin{proof}
  \uses{thm:T4, lem:T4-split}
  Lemma~\ref{lem:T4-split} and the triangle inequality bound each defect by
  $\tau^n+\sigma^n+\gamma^n$; the weights $L^{N-n}$ are non-negative.
\end{proof}

\begin{corollary}[T4, the contractive regime]\label{cor:T4-contractive}
  \lean{Atlas.master_bound_contractive}
  \leanok
  \uses{thm:T4}
  If moreover $0\le L<1$ and every one-step defect has norm at most $\delta$,
  then $\norm{u^N-u^{N,\star}}\le L^N\norm{u^0-u^{0,\star}}+\delta/(1-L)$ for
  every $N$.
\end{corollary}

\begin{proof}
  \uses{thm:T4}
  $\sum_{n=1}^{N}L^{N-n}=(1-L^N)/(1-L)\le1/(1-L)$.
\end{proof}

\begin{corollary}[T4, the non-expansive regime]\label{cor:T4-nonexpansive}
  \lean{Atlas.master_bound_nonexpansive}
  \leanok
  \uses{thm:T4}
  If $L=1$ and every one-step defect has norm at most $\delta$, then
  $\norm{u^N-u^{N,\star}}\le\norm{u^0-u^{0,\star}}+N\delta$.
\end{corollary}

\begin{proof}
  \uses{thm:T4}
  Every weight is $1$.
\end{proof}

\begin{corollary}[T4, the expansive regime]\label{cor:T4-expansive}
  \lean{Atlas.master_bound_expansive}
  \leanok
  \uses{thm:T4}
  If $L>1$ and every one-step defect has norm at most $\delta$, then
  $\norm{u^N-u^{N,\star}}\le L^N\norm{u^0-u^{0,\star}}
  +\dfrac{L^N-1}{L-1}\,\delta$.
\end{corollary}

\begin{proof}
  \uses{thm:T4}
  The geometric sum.
\end{proof}

\begin{proposition}[The regimes are attained]\label{prop:T4-attained}
  \lean{Atlas.master_bound_attained}
  \leanok
  For every $L,\delta\in\R$ and every $N$ there are real sequences with
  $u^{n+1}=L\,u^n$, a one-step defect equal to $\delta$ at every step and no
  initial error, whose error after $N$ steps is exactly
  $\sum_{i<N}L^i\,\delta$.
\end{proposition}

\begin{proof}
  $u^n=0$ and $u^{n+1,\star}=L\,u^{n,\star}-\delta$ with $u^{0,\star}=0$.
\end{proof}

\section{T4e: the transmission term}

\paragraph{In plain words.}
The term $\sigma$ comes from solving the interface problem with an approximate
operator $\tilde\Lambda$ in place of the exact one $\Lambda$.  Iterating does
not reduce it.  It is bounded by how wrong the operator is, amplified by
$1/\beta$, where $\beta$ is the smallest stretch of $\tilde\Lambda$ (its
smallest singular value).  An ill-conditioned interface makes the same operator
error cost more.

\begin{lemma}[T4e, the trace perturbation]\label{lem:T4e-trace}
  \lean{Atlas.trace_perturbation}
  \leanok
  Let $M,F$ be real normed spaces and $\Lambda,\tilde\Lambda:M\to F$ bounded and
  linear, with $\beta\norm{x}\le\norm{\tilde\Lambda x}$ for all $x$ and some
  $\beta>0$.  If $\Lambda\lambda^\star=\chi$ and
  $\tilde\Lambda\lambda^\dagger=\chi$, then
  \[
    \norm{\lambda^\dagger-\lambda^\star}
      \le\frac{\norm{\Lambda-\tilde\Lambda}\,\norm{\lambda^\star}}{\beta}.
  \]
\end{lemma}

\begin{proof}
  $\tilde\Lambda(\lambda^\dagger-\lambda^\star)=\chi-\tilde\Lambda\lambda^\star
  =(\Lambda-\tilde\Lambda)\lambda^\star$, so
  $\beta\norm{\lambda^\dagger-\lambda^\star}
  \le\norm{\Lambda-\tilde\Lambda}\norm{\lambda^\star}$.
\end{proof}

\begin{corollary}[T4e, the transmission term]\label{cor:T4e-sigma}
  \lean{Atlas.transmission_bound}
  \leanok
  \uses{lem:T4e-trace}
  If moreover $\norm{\Phi_\mu(v)-\Phi_{\mu'}(v)}\le C_\mu\norm{\mu-\mu'}$ for all
  data $\mu,\mu'$, then
  \[
    \norm{\Phi_{\lambda^\dagger}(v)-\Phi_{\lambda^\star}(v)}
      \le\frac{C_\mu}{\beta}\,\norm{\Lambda-\tilde\Lambda}\,\norm{\lambda^\star}.
  \]
\end{corollary}

\begin{proof}
  \uses{lem:T4e-trace}
  Compose the Lipschitz bound with Lemma~\ref{lem:T4e-trace}.
\end{proof}

\chapter{Decomposition does not change the answer}

A linear discrete problem $Au=f$ is cut into pieces.  Each piece is solved
exactly, given data on the cut.  The question of this chapter: when the pieces
have been made to agree, is the assembled field the solution of the undivided
problem?

For a direct interface solve and for Dirichlet--Neumann the answer is yes, with
no further hypothesis.  For restricted additive Schwarz it is yes under one
hypothesis, and false without it.

\section{T3a: a direct Schur interface solve}

\paragraph{In plain words.}
Split the unknowns into those inside the pieces and those on the interface.
Substructuring solves a smaller system for the interface unknowns alone, then
each piece solves for its interior with that interface data.  The result is
the undivided solution: nothing is lost and nothing is added by the cut.

\begin{theorem}[T3a, direct Schur solve]\label{thm:T3a}
  \lean{Atlas.schur_iff}
  \leanok
  Let $V_I,V_\Gamma$ be vector spaces over a field, and
  $A_{II}:V_I\to V_I$, $A_{I\Gamma}:V_\Gamma\to V_I$,
  $A_{\Gamma I}:V_I\to V_\Gamma$, $A_{\Gamma\Gamma}:V_\Gamma\to V_\Gamma$ linear,
  with $A_{II}$ invertible.  For all $f_I,u_I\in V_I$ and
  $f_\Gamma,u_\Gamma\in V_\Gamma$,
  \[
    \begin{cases}A_{II}u_I+A_{I\Gamma}u_\Gamma=f_I\\
    A_{\Gamma I}u_I+A_{\Gamma\Gamma}u_\Gamma=f_\Gamma\end{cases}
    \iff
    \begin{cases}u_I=A_{II}^{-1}(f_I-A_{I\Gamma}u_\Gamma)\\
    \bigl(A_{\Gamma\Gamma}-A_{\Gamma I}A_{II}^{-1}A_{I\Gamma}\bigr)u_\Gamma
      =f_\Gamma-A_{\Gamma I}A_{II}^{-1}f_I .\end{cases}
  \]
\end{theorem}

\begin{proof}
  \leanok
  The first equation on the left is equivalent to the first on the right
  because $A_{II}$ is invertible.  Substituting it into the second equation on
  the left gives the second on the right, and conversely.
\end{proof}

\section{T3b: Dirichlet--Neumann}

\paragraph{In plain words.}
Two pieces meet along an interface and do not overlap.  One sweep: piece 1 is
solved with the interface values given; the flux it then sends through the
interface is handed to piece 2, which returns its own interface values; the
interface values move toward those by a relaxation factor $\theta$.  If a sweep
returns the interface values it was given, the two pieces' solutions side by
side solve the undivided problem.  The solves must be exact.  Convergence of
the iteration is not assumed.

\begin{definition}[Two pieces in block form]\label{def:two-pieces}
  \lean{Atlas.TwoPieces, Atlas.TwoPieces.Solves, Atlas.TwoPieces.dnStep}
  \leanok
  The unknowns are $(u_1,u_2,\lambda)\in V_1\times V_2\times V_\Gamma$.  The
  undivided system is
  \begin{align*}
    A_{11}u_1+A_{1\Gamma}\lambda&=f_1,\\
    A_{22}u_2+A_{2\Gamma}\lambda&=f_2,\\
    A_{\Gamma1}u_1+A_{\Gamma2}u_2
      +\bigl(A^{(1)}_{\Gamma\Gamma}+A^{(2)}_{\Gamma\Gamma}\bigr)\lambda
      &=f^{(1)}_\Gamma+f^{(2)}_\Gamma ,
  \end{align*}
  where each piece contributes its own share of the interface block and of the
  interface right-hand side.  The flux piece 1 sends is
  $q(u_1,\lambda)=A_{\Gamma1}u_1+A^{(1)}_{\Gamma\Gamma}\lambda-f^{(1)}_\Gamma$.
  A Dirichlet solve $D$ is exact if $A_{11}D(\lambda)+A_{1\Gamma}\lambda=f_1$
  for all $\lambda$.  A Neumann solve $N$, returning a pair
  $(u_2,\tilde\lambda)$, is exact if
  $A_{22}u_2+A_{2\Gamma}\tilde\lambda=f_2$ and
  $A_{\Gamma2}u_2+A^{(2)}_{\Gamma\Gamma}\tilde\lambda=f^{(2)}_\Gamma-q$ for all
  $q$.  One sweep is
  $\lambda\mapsto\lambda+\theta\,(\tilde\lambda-\lambda)$ with
  $(u_2,\tilde\lambda)=N\bigl(q(D(\lambda),\lambda)\bigr)$.
\end{definition}

\begin{theorem}[T3b, Dirichlet--Neumann, block form]\label{thm:T3b}
  \lean{Atlas.TwoPieces.dn_fixedPoint_solves}
  \leanok
  \uses{def:two-pieces}
  With exact solves and $\theta\ne0$, if $\lambda$ is a fixed point of the
  sweep then $\bigl(D(\lambda),u_2,\lambda\bigr)$ solves the undivided system,
  where $u_2$ is the interior the Neumann solve returned.
\end{theorem}

\begin{proof}
  \uses{def:two-pieces}
  \leanok
  $\theta\ne0$ gives $\tilde\lambda=\lambda$.  The two interior equations are
  the exactness of the solves.  The Neumann solve's interface equation,
  $A_{\Gamma2}u_2+A^{(2)}_{\Gamma\Gamma}\lambda=f^{(2)}_\Gamma-q$, with $q$
  written out, is the interface equation of the undivided system.
\end{proof}

\begin{theorem}[T3b, converse]\label{thm:T3b-converse}
  \lean{Atlas.TwoPieces.dn_solution_isFixedPt}
  \leanok
  \uses{def:two-pieces}
  If each piece's equations determine its solution and the solves return it,
  then for every solution $(u_1,u_2,\lambda)$ of the undivided system and every
  $\theta$, $\lambda$ is a fixed point of the sweep.
\end{theorem}

\begin{proof}
  \uses{def:two-pieces}
  \leanok
  The Dirichlet solve returns $u_1$.  The pair $(u_2,\lambda)$ satisfies the
  Neumann piece's equations with the flux $q(u_1,\lambda)$, so the Neumann solve
  returns it, and the update is zero.
\end{proof}

\begin{theorem}[T3b, the face form the workbench implements]\label{thm:T3b-face}
  \lean{Atlas.FacePair.solves_of_fixedPoint}
  \leanok
  Two sets of finite-volume cells, $D$ and $N$, meet along faces.  $K_D,K_N$ are
  each set's operator with the cut faces left out; $P_D,P_N$ read the cell value
  beside each cut face; $Q_D,Q_N$ put a per-face flow into those cells'
  equations; $g_D$ is the $D$ side's half-cell conductance, $r_D=g_D^{-1}$ and
  $r_N$ the two half-cell resistances, and $h=(r_D+r_N)^{-1}$ the undivided
  grid's face conductance (the harmonic mean).  Suppose
  \begin{enumerate}
    \item $K_Du_D+Q_Dg_DP_Du_D=b_D+Q_Dg_D\lambda$ (the $D$ set solved with face
      values $\lambda$);
    \item $q=g_D(P_Du_D-\lambda)$ (the flow it sends);
    \item $K_Nu_N=b_N+Q_Nq$ (the $N$ set solved with that flow);
    \item $P_Nu_N+r_Nq=\lambda$ (the $N$ set's face values reproduce
      $\lambda$).
  \end{enumerate}
  Then $q=h\,(P_Du_D-P_Nu_N)$, and $(u_D,u_N)$ solves the undivided system
  \[
    K_Du_D+Q_Dh(P_Du_D-P_Nu_N)=b_D,\qquad
    K_Nu_N-Q_Nh(P_Du_D-P_Nu_N)=b_N .
  \]
\end{theorem}

\begin{proof}
  \leanok
  From 2, $r_Dq=P_Du_D-\lambda$; from 4, $r_Nq=\lambda-P_Nu_N$.  Adding,
  $(r_D+r_N)q=P_Du_D-P_Nu_N$, so $q=h(P_Du_D-P_Nu_N)$.  Then 1 reads
  $K_Du_D+Q_Dq=b_D$ and 3 reads $K_Nu_N-Q_Nq=b_N$.
\end{proof}

\section{T3c: restricted additive Schwarz}

\paragraph{In plain words.}
The domain is covered by overlapping windows.  One sweep solves every window
exactly, with the current iterate held outside it, and blends the windows'
solutions with weights that sum to one.

\begin{itemize}
  \item The undivided solution is always a fixed point of the sweep.
  \item If, at $u$, every window's own solution agrees with $u$ on the whole
    window, then $u$ is the undivided solution.
  \item If only the \emph{blend} is unchanged, more is needed.  The sweep is
    $u\mapsto u+M(f-Au)$, with $M$ the blended sum of the window solves, and a
    fixed point is the undivided solution exactly when $M$ sends no non-zero
    residual to zero.  That holds whenever the sweep converges from every
    start.
  \item Without that hypothesis the claim is false
    (Theorem~\ref{thm:T3c-counterexample}).
\end{itemize}

\begin{definition}[A Schwarz sweep]\label{def:schwarz}
  \lean{Atlas.Schwarz, Atlas.Schwarz.window, Atlas.Schwarz.sweep, Atlas.Schwarz.precond}
  \leanok
  $A:V\to V$ is linear.  For each window $i$ of a finite family:
  $R_i:V\to W_i$ reads the window's unknowns; $E_i:W_i\to V$ pads with zeros;
  $E^\chi_i:W_i\to V$ multiplies by the window's weights, then pads;
  $B_i:W_i\to W_i$ is the inverse of the window's operator $R_iAE_i$.  The
  weights are a partition of unity: $\sum_iE^\chi_iR_i=I$.  The window's
  solution, the sweep and the preconditioner are
  \[
    w_i(f,u)=B_i\bigl(R_if-R_iA(u-E_iR_iu)\bigr),\qquad
    G_f(u)=\sum_iE^\chi_i\,w_i(f,u),\qquad
    M=\sum_iE^\chi_iB_iR_i .
  \]
\end{definition}

\begin{lemma}[The sweep in residual form]\label{lem:T3c-sweep}
  \lean{Atlas.Schwarz.sweep_eq}
  \leanok
  \uses{def:schwarz}
  $G_f(u)=u+M(f-Au)$.
\end{lemma}

\begin{proof}
  \uses{def:schwarz}
  \leanok
  $w_i(f,u)=B_iR_i(f-Au)+B_i(R_iAE_i)R_iu=B_iR_i(f-Au)+R_iu$.  Blend, and use
  the partition of unity.
\end{proof}

\begin{theorem}[T3c, consistency]\label{thm:T3c-consistency}
  \lean{Atlas.Schwarz.solution_isFixedPt}
  \leanok
  \uses{lem:T3c-sweep}
  If $Au=f$ then $G_f(u)=u$.
\end{theorem}

\begin{proof}
  \uses{lem:T3c-sweep}
  \leanok
  $M0=0$.
\end{proof}

\begin{lemma}[What a fixed point is]\label{lem:T3c-iff}
  \lean{Atlas.Schwarz.isFixedPt_iff}
  \leanok
  \uses{lem:T3c-sweep}
  $G_f(u)=u\iff M(f-Au)=0$.
\end{lemma}

\begin{proof}
  \uses{lem:T3c-sweep}
  \leanok
  Lemma~\ref{lem:T3c-sweep}.
\end{proof}

\begin{theorem}[T3c, agreement]\label{thm:T3c-agree}
  \lean{Atlas.Schwarz.solves_of_windows_agree}
  \leanok
  \uses{def:schwarz}
  If $w_i(f,u)=R_iu$ for every window $i$, then $Au=f$.
\end{theorem}

\begin{proof}
  \uses{def:schwarz}
  \leanok
  As in Lemma~\ref{lem:T3c-sweep}, $w_i(f,u)=R_iu+B_iR_i(f-Au)$, so
  $B_iR_i(f-Au)=0$ and, $B_i$ being invertible, $R_i(f-Au)=0$ for every $i$.
  By the partition of unity $f-Au=\sum_iE^\chi_iR_i(f-Au)=0$.
\end{proof}

\begin{theorem}[T3c, a stationary blend]\label{thm:T3c-precond}
  \lean{Atlas.Schwarz.fixedPt_solves_of_precond}
  \leanok
  \uses{lem:T3c-iff}
  If $Mr=0$ implies $r=0$, then every fixed point of $G_f$ solves $Au=f$.
\end{theorem}

\begin{proof}
  \uses{lem:T3c-iff}
  \leanok
  Lemma~\ref{lem:T3c-iff}.
\end{proof}

\begin{corollary}[The same hypothesis, read off the sweep]\label{cor:T3c-homogeneous}
  \lean{Atlas.Schwarz.fixedPt_solves}
  \leanok
  \uses{thm:T3c-precond, lem:T3c-sweep}
  If $A$ is onto and the sweep with zero right-hand side has no fixed point but
  $0$, then for every $f$ every fixed point of $G_f$ solves $Au=f$.
\end{corollary}

\begin{proof}
  \uses{thm:T3c-precond, lem:T3c-sweep}
  \leanok
  Let $Mr=0$ and write $r=Ae$.  Then $G_0(-e)=-e+M(Ae)=-e$, so $e=0$ and
  $r=0$.
\end{proof}

\begin{theorem}[T3c, a convergent sweep has one fixed point]\label{thm:T3c-unique}
  \lean{Atlas.Schwarz.fixedPt_eq_solution}
  \leanok
  \uses{lem:T3c-sweep}
  Let $V$ be a real normed space, and suppose $G_0^{\,k}(e)\to0$ for every
  $e\in V$.  If $Au^\star=f$ and $G_f(u)=u$, then $u=u^\star$.
\end{theorem}

\begin{proof}
  \uses{lem:T3c-sweep}
  \leanok
  $G_f(u)-G_f(u^\star)=G_0(u-u^\star)$, so $e=u-u^\star$ is a fixed point of
  $G_0$.  Then $e=G_0^{\,k}(e)\to0$.
\end{proof}

\begin{theorem}[T3c, a convergent sweep converges to the solution]\label{thm:T3c-tendsto}
  \lean{Atlas.Schwarz.tendsto_solution}
  \leanok
  \uses{lem:T3c-sweep, thm:T3c-consistency}
  Under the same hypothesis, if $Au^\star=f$ then $G_f^{\,k}(u_0)\to u^\star$
  for every $u_0\in V$.
\end{theorem}

\begin{proof}
  \uses{lem:T3c-sweep, thm:T3c-consistency}
  \leanok
  $G_f^{\,k}(u_0)-u^\star=G_0^{\,k}(u_0-u^\star)$.
\end{proof}

\begin{theorem}[T3c, the counterexample]\label{thm:T3c-counterexample}
  \lean{Atlas.Schwarz.exists_spurious_fixedPt}
  \leanok
  \uses{def:schwarz}
  There is a Schwarz sweep over $\Q$ with exact window solves, a partition of
  unity and an invertible $A$, and a pair $f,u$ with $G_f(u)=u$ and $Au\ne f$.
\end{theorem}

\begin{proof}
  \uses{def:schwarz}
  \leanok
  Three unknowns, windows $\{1,2\}$ and $\{2,3\}$, the shared unknown given
  entirely to the first window:
  \[
    A=\begin{pmatrix}1&1&0\\1&0&1\\0&1&1\end{pmatrix},\qquad f=0,\qquad
    u=(1,-1,-1).
  \]
  $\det A=-2$ and both window matrices have determinant $-1$.  Window 1 returns
  $(1,-1)$, which is $u$ on $\{1,2\}$.  Window 2 returns $(1,-1)$ on $\{2,3\}$:
  it disagrees with $u$ at the shared unknown, and the blend discards that
  value.  So $G_0(u)=u$, while $Au=(0,0,-2)$.
\end{proof}

\begin{remark}[Positive definite matrices]
  What fails in the counterexample is that $A$ has a zero on its diagonal at
  the shared unknown: the problem restricted to the overlap is singular.  For
  \emph{two} windows and a symmetric positive definite $A$ this cannot happen,
  and every fixed point is the solution (a paper proof, not yet formalised).
  With \emph{three} windows positive definiteness is not enough: the project's
  script \texttt{scripts/lean\_t3\_counterexample.py} exhibits, in exact
  rational arithmetic, a symmetric positive definite $6\times6$ matrix and three
  windows with a spurious fixed point.  On that example the sweep does not
  converge, so Theorem~\ref{thm:T3c-unique} is not contradicted.  For
  M-matrices the sweep always converges (Frommer and Szyld, \emph{SIAM J.
  Numer. Anal.} 39, 2001).
\end{remark}

\section{T3d: what stopping on the update leaves}

\paragraph{In plain words.}
A sweep is stopped when its update is small.  The update is not the error.  If
the zero-right-hand-side sweep shrinks a norm by $\rho<1$ every time, the error
is at most the update divided by $1-\rho$.  For a slowly converging sweep,
$\rho$ near one, that is many times the tolerance.

\begin{theorem}[T3d, stopping on the update]\label{thm:T3d}
  \lean{Atlas.Schwarz.error_le_update}
  \leanok
  \uses{thm:T7, lem:T3c-sweep}
  Let $V$ be a real normed space and $\norm{G_0(e)}\le\rho\norm{e}$ for all
  $e$, with $\rho<1$.  If $Au^\star=f$ then for every $u\in V$
  \[
    \norm{u-u^\star}\le\frac{\norm{G_f(u)-u}}{1-\rho}.
  \]
\end{theorem}

\begin{proof}
  \uses{thm:T7, lem:T3c-sweep}
  \leanok
  $G_f(u)-G_f(u^\star)=G_0(u-u^\star)$ and $G_f(u^\star)=u^\star$, so $G_f$
  contracts the pair $(u,u^\star)$ by $\rho$; apply Theorem~\ref{thm:T7}.
\end{proof}

\chapter{Defect correction: the rate, never the answer}

Every other slot for a learned expert asks it to be part of the answer, and so
owes it a certificate.  Defect correction puts it somewhere else.  $\Phi$ is the
classical step, whose settled states are its fixed points.  $\Psi$ is any other
map on the same states: a learned operator, a coarse solver, anything cheap.
With $G_\Phi=I-\Phi$ and $G_\Psi=I-\Psi$, the next iterate $w'$ is defined by
\[
  G_\Psi(w')=G_\Psi(w)-G_\Phi(w),\qquad\text{that is,}\qquad
  w'-\Psi(w')=\Phi(w)-\Psi(w)
\]
(Stetter, \emph{Numer. Math.} 29, 1978).  The cheap map enters only through a
difference, and the classical map supplies the defect.

\section{T5: Theorem 1 and Corollary 1}

\paragraph{In plain words.}
If the iterates settle anywhere, they settle on a fixed point of the classical
map.  Nothing about the accuracy of the cheap map enters, only that both maps
are continuous at the limit.  A wrong cheap map can cost classical calls; it
cannot change what is returned.  And if the cheap map is constant, the
iteration is the classical march itself.

\begin{definition}[A defect-correction step]\label{def:dec-step}
  \lean{Atlas.DefectCorrection.IsStep, Atlas.DefectCorrection.innerMarch}
  \leanok
  Let $E$ be a normed space and $\Phi,\Psi:E\to E$.  The point $w'$ is a
  successor of $w$ if $w'-\Psi(w')=(w-\Psi(w))-(w-\Phi(w))$.  The \emph{inner
  march} that solves for it is $x^{(0)}=\Phi(w)$,
  $x^{(m+1)}=\Phi(w)+[\Psi(x^{(m)})-\Psi(w)]$.
\end{definition}

\begin{theorem}[T5, Theorem 1: consistency, whatever $\Psi$ is]\label{thm:T5}
  \lean{Atlas.DefectCorrection.limit_isFixedPt}
  \leanok
  \uses{def:dec-step}
  Let $w_{k+1}$ be a successor of $w_k$ for every $k$, and $w_k\to\bar w$.  If
  $\Phi$ and $\Psi$ are continuous at $\bar w$, then $\Phi(\bar w)=\bar w$.
\end{theorem}

\begin{proof}
  \uses{def:dec-step}
  $w_k-\Phi(w_k)=\bigl(w_k-\Psi(w_k)\bigr)-\bigl(w_{k+1}-\Psi(w_{k+1})\bigr)$.
  The right side tends to $(\bar w-\Psi(\bar w))-(\bar w-\Psi(\bar w))=0$ and the
  left to $\bar w-\Phi(\bar w)$.
\end{proof}

\begin{theorem}[T5, Theorem 1 with inexact inner solves]\label{thm:T5-inexact}
  \lean{Atlas.DefectCorrection.limit_isFixedPt_of_inexact}
  \leanok
  \uses{def:dec-step}
  Suppose instead that each step misses the defining equation by at most
  $\varepsilon_k$,
  \[
    \norm{\bigl(w_{k+1}-\Psi(w_{k+1})\bigr)
      -\bigl((w_k-\Psi(w_k))-(w_k-\Phi(w_k))\bigr)}\le\varepsilon_k ,
  \]
  with $\varepsilon_k\to0$.  If $w_k\to\bar w$ and $\Phi$, $\Psi$ are continuous
  at $\bar w$, then $\Phi(\bar w)=\bar w$.
\end{theorem}

\begin{proof}
  \uses{def:dec-step}
  As for Theorem~\ref{thm:T5}: the left side of the same identity now differs
  from the right by at most $\varepsilon_k\to0$.
\end{proof}

\begin{corollary}[T5, Corollary 1: the null element]\label{cor:T5-const}
  \lean{Atlas.DefectCorrection.step_of_const}
  \leanok
  \uses{def:dec-step}
  If $\Psi$ is constant and $w'$ is a successor of $w$, then $w'=\Phi(w)$.
\end{corollary}

\begin{proof}
  \uses{def:dec-step}
  $w'-c=(w-c)-(w-\Phi(w))$.
\end{proof}

\begin{lemma}[The inner march solves for the successor]\label{lem:T5-inner}
  \lean{Atlas.DefectCorrection.isStep_iff}
  \leanok
  \uses{def:dec-step}
  $x$ is a successor of $w$ if and only if $x=\Phi(w)+[\Psi(x)-\Psi(w)]$.
\end{lemma}

\begin{proof}
  \uses{def:dec-step}
  Rearrangement.
\end{proof}

\begin{corollary}[T5, Corollary 1 for the inner march]\label{cor:T5-inner-const}
  \lean{Atlas.DefectCorrection.innerMarch_of_const}
  \leanok
  \uses{def:dec-step}
  If $\Psi$ is constant, $x^{(m)}=\Phi(w)$ for every $m$.
\end{corollary}

\begin{proof}
  \uses{def:dec-step}
  Induction on $m$; the bracket is zero.
\end{proof}

\section{T5: Corollary 2, shrinking toward the null element}

\paragraph{In plain words.}
The iteration's rate is set by $J_\Psi^{-1}J_\Phi$, with $J_\Psi=I-D\Psi$ the
derivative of $G_\Psi$ at the settled state.  A cheap map that is nearly
singular on some mode makes the iteration diverge.  Replace $\Psi$ by
$(1-\alpha)\Psi$.  Its $J$ is $\alpha I+(1-\alpha)J_\Psi$, whose eigenvalues are
those of $J_\Psi$ moved toward $1$; wherever $J_\Psi$'s eigenvalue has
non-negative real part, the new one has real part at least $\alpha$.

\begin{lemma}[The shrunk Jacobian]\label{lem:T5-shrink-deriv}
  \lean{Atlas.DefectCorrection.shrink_hasFDerivAt}
  \leanok
  If $\Psi$ has Fr\'echet derivative $D$ at $w$, then
  $x\mapsto x-(1-\alpha)\Psi(x)$ has derivative
  $\alpha I+(1-\alpha)(I-D)$ at $w$.
\end{lemma}

\begin{proof}
  Linearity of the derivative, and
  $I-(1-\alpha)D=\alpha I+(1-\alpha)(I-D)$.
\end{proof}

\begin{theorem}[T5, Corollary 2: the eigenvalues]\label{thm:T5-shrink}
  \lean{Atlas.DefectCorrection.shrink_hasEigenvalue_iff}
  \leanok
  \uses{lem:T5-shrink-eig}
  Let $J$ be a linear map on a vector space over a field, and $\alpha\ne1$.
  Then $\nu$ is an eigenvalue of $\alpha I+(1-\alpha)J$ if and only if
  $\nu=\alpha+(1-\alpha)\mu$ for an eigenvalue $\mu$ of $J$.
\end{theorem}

\begin{proof}
  \uses{lem:T5-shrink-eig}
  Lemma~\ref{lem:T5-shrink-eig} gives one direction.  For the other,
  $(\alpha I+(1-\alpha)J)v=\nu v$ gives $Jv=\frac{\nu-\alpha}{1-\alpha}v$.
\end{proof}

\begin{lemma}[Eigenvalues move with the shrink]\label{lem:T5-shrink-eig}
  \lean{Atlas.DefectCorrection.shrink_hasEigenvalue}
  \leanok
  If $Jv=\mu v$ then $(\alpha I+(1-\alpha)J)v=(\alpha+(1-\alpha)\mu)v$.
\end{lemma}

\begin{proof}
  Direct.
\end{proof}

\begin{lemma}[The real part]\label{lem:T5-shrink-re}
  \lean{Atlas.DefectCorrection.shrink_re}
  \leanok
  For $0\le\alpha\le1$ and complex $\mu$ with $\operatorname{Re}\mu\ge0$,
  $\operatorname{Re}\bigl(\alpha+(1-\alpha)\mu\bigr)\ge\alpha$.
\end{lemma}

\begin{proof}
  $\operatorname{Re}(\alpha+(1-\alpha)\mu)=\alpha+(1-\alpha)\operatorname{Re}\mu$.
\end{proof}

\chapter{Guarantees the composition layer supplies}

Two properties that no single expert has to possess, because the way experts
are called and connected creates them: an exact symmetry, and passivity.

\section{T10: symmetry averaging}

\paragraph{In plain words.}
A frozen expert does not respect a symmetry.  Call it once for every element of
the symmetry group, undo each transformation, and average.  The averaged map
respects the symmetry exactly, whatever the expert is: nonlinear, learned,
discontinuous.  The proof uses only that multiplying every group element by a
fixed one runs through the whole group again, which is why a set of operations
that is not a group does not work.

\begin{definition}[The group average]\label{def:average}
  \lean{Atlas.average}
  \leanok
  Let $G$ be a finite group acting on a set $V$ and acting linearly on a vector
  space $W$.  For any map $E:V\to W$,
  \[
    \tilde E(u)=\frac1{|G|}\sum_{g\in G}g^{-1}\,E(g\,u).
  \]
\end{definition}

\begin{theorem}[T10, symmetry averaging]\label{thm:T10}
  \lean{Atlas.average_equivariant}
  \leanok
  \uses{def:average}
  For every map $E$, every $h\in G$ and every $u\in V$,
  $\tilde E(h\,u)=h\,\tilde E(u)$.
\end{theorem}

\begin{proof}
  \uses{def:average}
  \leanok
  $\tilde E(hu)=\frac1{|G|}\sum_gg^{-1}E(ghu)$.  Substitute $g'=gh$, which runs
  over $G$ as $g$ does; then $g^{-1}=h\,g'^{-1}$, and $h$ comes out of the sum
  because the action on $W$ is linear.
\end{proof}

\begin{proposition}[Averaging leaves a symmetric map alone]\label{prop:T10-idem}
  \lean{Atlas.average_of_equivariant}
  \leanok
  \uses{def:average}
  If $E(gu)=gE(u)$ for all $g,u$, and $|G|$ is invertible in the scalars, then
  $\tilde E=E$.
\end{proposition}

\begin{proof}
  \uses{def:average}
  \leanok
  Every term of the sum is $E(u)$.
\end{proof}

\begin{example}[A set that is not a group]\label{ex:T10-nongroup}
  \lean{Atlas.average_over_non_group_not_equivariant}
  \leanok
  Let $M_x(x,y)=(-x,y)$ and $M_y(x,y)=(x,-y)$ on $\R^2$.  The set
  $\{\mathrm{id},M_x,M_y\}$ is not a group: $M_xM_y$ is missing.  The same
  average over this set is not $M_x$-equivariant: for $E(x,y)=(x+y,0)$ and
  $u=(0,1)$, the average at $M_xu=u$ is $(-\tfrac13,0)$, and $M_x$ of the
  average at $u$ is $(\tfrac13,0)$.
\end{example}

\begin{proof}
  \leanok
  $E(u)=(1,0)$, $M_xE(M_xu)=M_x(1,0)=(-1,0)$ and $M_yE(M_yu)=M_y(-1,0)=(-1,0)$,
  so the average at $u=M_xu$ is $\tfrac13(-1,0)$, and $M_x$ of it is
  $(\tfrac13,0)\ne(-\tfrac13,0)$.
\end{proof}

\section{T11: passive agents joined by the port rule}

\paragraph{In plain words.}
A port carries an effort and a flow whose inner product is power.  The port
algebra joins two ports by one rule: equal efforts, opposite flows.  Then the
power entering one port is exactly the power leaving its partner, so over all
connected ports the powers cancel, whatever the graph.  If every agent is
passive (its stored energy grows by no more than what entered through its
ports), the whole system is passive with respect to its open ports.  And if
every agent is \emph{incrementally} passive, one coupled step does not increase
the distance between two runs: the constant $L\le1$ of the master bound.

\begin{lemma}[The connection rule conserves power]\label{lem:T11-rule}
  \lean{Atlas.port_rule_power}
  \leanok
  In a real inner-product space, if $e_B=e_A$ and $f_B=-f_A$ then
  $\langle e_B,f_B\rangle=-\langle e_A,f_A\rangle$.
\end{lemma}

\begin{proof}
  \leanok
  Linearity of the inner product in its second argument.
\end{proof}

\begin{lemma}[The junction]\label{lem:T11-junction}
  \lean{Atlas.junction_power_eq_zero}
  \leanok
  Let $P$ be a finite set of ports, $\pi:P\to P$ a bijection pairing each port
  with its partner, and $p:P\to\R$ with $p(\pi(x))=-p(x)$ for all $x$.  Then
  $\sum_{x\in P}p(x)=0$.
\end{lemma}

\begin{proof}
  \leanok
  $\sum_xp(x)=\sum_xp(\pi(x))=-\sum_xp(x)$.
\end{proof}

\begin{theorem}[T11, passivity composes]\label{thm:T11}
  \lean{Atlas.interconnection_passive}
  \leanok
  \uses{lem:T11-junction}
  Let each port belong to one of finitely many agents.  Agent $i$ stores $H_i$
  before a step and $H_i'$ after it, receives the power of its own connected
  ports and $\mathrm{ext}_i$ through its open ports, and is passive:
  $H_i'\le H_i+\sum_{x\text{ of }i}p(x)+\mathrm{ext}_i$.  If the connected ports
  obey $p(\pi(x))=-p(x)$, then
  \[
    \sum_iH_i'\le\sum_iH_i+\sum_i\mathrm{ext}_i .
  \]
\end{theorem}

\begin{proof}
  \uses{lem:T11-junction}
  \leanok
  Sum the agents' inequalities.  The port powers, summed over agents, are the
  sum over all ports, which is zero by Lemma~\ref{lem:T11-junction}.
\end{proof}

\begin{corollary}[T11, incremental passivity gives a non-expansive step]\label{cor:T11-nonexpansive}
  \lean{Atlas.interconnection_nonexpansive}
  \leanok
  \uses{thm:T11}
  Let $u_i,v_i$ be agent $i$'s states in two runs before a step and $u_i',v_i'$
  after it, and let $\Delta p(x)$ be the power pairing of the differences of
  the two runs' efforts and flows at port $x$.  If
  $\norm{u_i'-v_i'}^2\le\norm{u_i-v_i}^2+2\sum_{x\text{ of }i}\Delta p(x)$ for
  every agent and $\Delta p(\pi(x))=-\Delta p(x)$, then
  \[
    \sum_i\norm{u_i'-v_i'}^2\le\sum_i\norm{u_i-v_i}^2 .
  \]
\end{corollary}

\begin{proof}
  \uses{thm:T11}
  \leanok
  Theorem~\ref{thm:T11} with $H_i=\norm{u_i-v_i}^2$, the powers $2\Delta p$, and
  no external term.
\end{proof}

\chapter{Tier 0: the learned superelement}

At tier 0 a learned expert is not asked for a port matrix.  It returns
\emph{fields}: for each pattern of boundary values, an interior field to go
with it.  The host computes the port matrix itself, as the energy Gram matrix
of those fields, assembles one interface system, and solves it once.  There is
no iteration.

The statements of this chapter hold for \emph{every} output of the network,
however wrong.  They are the tier-0 guarantee of the architecture document
(its Theorems 5 and 6, Proposition 7 and Corollary 11): the learned system is
never indefinite, it is solvable whenever the classical one is, its answer is
the best its modes can represent, and its sensitivity is governed by the
classical problem's own constant.

Everything is finite-dimensional linear algebra.  The statements are written
for real vector spaces and bilinear forms, so that nothing depends on a choice
of basis for the fields; with coordinates they are the matrix statements of the
document.

\section{T25 (i) to (iii): the Gram port matrix}

\paragraph{In plain words.}
A piece of the domain has unknowns inside it and unknowns on its boundary
faces, and an energy that is never negative.  A port mode is a pattern of
boundary values; its exact constraint mode is the interior field the physics
pairs with it, and the exact port matrix is the energy Gram matrix of those
pairs.  Put any other interior fields in place of the exact ones.  Then the
Gram matrix is still symmetric and positive semidefinite; it exceeds the exact
one by exactly the Gram matrix of the field errors in the interior energy; and
after assembly over all pieces the learned interface matrix dominates the exact
one, so it inherits every lower bound the exact one has.  An error in the
fields can make a piece stiffer, never softer, and never indefinite.

\begin{definition}[The energy of a piece]\label{def:piece}
  \lean{Atlas.Piece, Atlas.Piece.M, Atlas.Piece.IsExactMode}
  \leanok
  Let $U$ (cell values) and $F$ (face values) be real vector spaces.  A
  \emph{piece} is three bilinear forms, written with matrices as
  $a(u,v)=u^\top A_Iv$ on $U$, $b(u,\mu)=u^\top B\mu$ on $U\times F$ and
  $d(\lambda,\mu)=\lambda^\top D\mu$ on $F$, with $a$ and $d$ symmetric and
  \[
    a(u,u)-2\,b(u,\lambda)+d(\lambda,\lambda)\ \ge\ 0
    \qquad\text{for all }u\in U,\ \lambda\in F .
  \]
  Its energy form on pairs is
  \[
    M\bigl((u,\lambda),(v,\mu)\bigr)=a(u,v)-b(u,\mu)-b(v,\lambda)+d(\lambda,\mu),
    \qquad M=\begin{pmatrix}A_I&-B\\-B^\top&D\end{pmatrix}.
  \]
  An interior field $h\in U$ is the \emph{exact constraint mode} of the
  boundary pattern $q\in F$ if $a(v,h)=b(v,q)$ for every $v\in U$, that is,
  $A_Ih=Bq$.
\end{definition}

\begin{definition}[The Gram port matrix]\label{def:gram}
  \lean{Atlas.Piece.gram}
  \leanok
  \uses{def:piece}
  For interior fields $h_k\in U$ and boundary patterns $q_k\in F$, indexed by a
  finite set, the Gram port matrix has entries
  \[
    \Lambda[h]_{kl}=M\bigl((h_k,q_k),(h_l,q_l)\bigr).
  \]
  With the exact constraint modes $H_k$ it is the exact port matrix
  $\Lambda=\Lambda[H]$; with a network's fields $\hat H_k$ it is the matrix
  $\tilde\Lambda=\Lambda[\hat H]$ the host computes.
\end{definition}

\begin{theorem}[T25 (i), positive semidefinite for every output]\label{thm:T25-i}
  \lean{Atlas.Piece.gram_posSemidef}
  \leanok
  \uses{def:gram}
  For any fields $h_k$ and patterns $q_k$, $\Lambda[h]$ is symmetric and
  $x^\top\Lambda[h]\,x\ge0$ for every $x$.
\end{theorem}

\begin{proof}
  \leanok
  $\Lambda[h]$ is symmetric because $M$ is.  With
  $w=\bigl(\sum_kx_kh_k,\ \sum_kx_kq_k\bigr)$, bilinearity gives
  $x^\top\Lambda[h]\,x=M(w,w)\ge0$.
\end{proof}

\begin{proposition}[T25 (i), the definite case]\label{prop:T25-i-def}
  \lean{Atlas.Piece.gram_posDef}
  \leanok
  \uses{def:gram}
  If $M(w,w)=0$ only for $w=0$, and the patterns $q_k$ are linearly
  independent, then $\Lambda[h]$ is positive definite for any fields $h_k$.
\end{proposition}

\begin{proof}
  \leanok
  \uses{thm:T25-i}
  If $x^\top\Lambda[h]\,x=0$ then $w=0$ in the proof of
  Theorem~\ref{thm:T25-i}; its face part is $\sum_kx_kq_k=0$, so $x=0$.
\end{proof}

\begin{theorem}[T25 (ii), the error identity]\label{thm:T25-ii}
  \lean{Atlas.Piece.gram_sub_gram, Atlas.Piece.quadratic_gram_sub}
  \leanok
  \uses{def:gram}
  Let $H_k$ be the exact constraint modes of the patterns $q_k$, let $h_k$ be
  any fields, and $E_k=h_k-H_k$.  Then
  \[
    \Lambda[h]_{kl}-\Lambda[H]_{kl}=a(E_k,E_l)\qquad\text{for all }k,l,
  \]
  and so, for every coordinate vector $x$,
  \[
    x^\top\bigl(\Lambda[h]-\Lambda[H]\bigr)x
      =a\Bigl(\sum_kx_kE_k,\ \sum_kx_kE_k\Bigr).
  \]
  In matrices: $\tilde\Lambda-\Lambda=E^\top A_IE$.
\end{theorem}

\begin{proof}
  \leanok
  Write $(h_k,q_k)=(H_k,q_k)+(E_k,0)$ and expand $M$ in both arguments.  The
  cross terms are $M\bigl((E_k,0),(H_l,q_l)\bigr)=a(E_k,H_l)-b(E_k,q_l)=0$ by
  the definition of $H_l$, and the same with $k$ and $l$ exchanged.  The last
  term is $M\bigl((E_k,0),(E_l,0)\bigr)=a(E_k,E_l)$.
\end{proof}

\begin{corollary}[T25 (ii), consequences]\label{cor:T25-ii}
  \lean{Atlas.Piece.gram_sub_gram_posSemidef, Atlas.Piece.trace_gram_sub,
    Atlas.Piece.quadratic_gram_sub_le}
  \leanok
  \uses{thm:T25-ii}
  With the notation of Theorem~\ref{thm:T25-ii}:
  \begin{enumerate}
    \item $\Lambda[h]-\Lambda[H]$ is positive semidefinite;
    \item $\operatorname{tr}\Lambda[h]-\operatorname{tr}\Lambda[H]=\sum_ka(E_k,E_k)$;
    \item if $s:U\to\R$ is any measure of size with $a(u,u)\le\alpha\,s(u)^2$
      for all $u$, $\alpha\ge0$, and
      $s\bigl(\sum_kx_kE_k\bigr)^2\le\varepsilon^2\,|x|^2$ for all $x$, then
      $x^\top\bigl(\Lambda[h]-\Lambda[H]\bigr)x\le\alpha\,\varepsilon^2\,|x|^2$.
      With Euclidean norms this is
      $\norm{\tilde\Lambda-\Lambda}_2\le\norm{A_I}_2\norm{E}_2^2$: the error of
      the port matrix is quadratic in the error of the fields.
  \end{enumerate}
\end{corollary}

\begin{proof}
  \leanok
  \uses{thm:T25-ii}
  The interior energy $a(u,u)=M\bigl((u,0),(u,0)\bigr)$ is non-negative, which
  gives 1 from the quadratic form of Theorem~\ref{thm:T25-ii}.  Item 2 is the
  diagonal of the identity.  Item 3 is the quadratic form with the two bounds.
\end{proof}

\begin{definition}[Assembly]\label{def:assemble}
  \lean{Atlas.assemble}
  \leanok
  For pieces $i$ in a finite set, port matrices $L_i$ and matrices $R_i$ that
  read piece $i$'s port coordinates off the global coordinate vector, the
  assembled interface matrix is $\sum_iR_i^\top L_iR_i$.
\end{definition}

\begin{theorem}[T25 (iii), the assembled matrix]\label{thm:T25-iii}
  \lean{Atlas.assemble_gram_sub_posSemidef, Atlas.assemble_gram_coercive,
    Atlas.assemble_gram_posDef, Atlas.assemble_gram_existsUnique}
  \leanok
  \uses{def:assemble, cor:T25-ii}
  Let $\mathsf S=\sum_iR_i^\top\Lambda_i[H_i]R_i$ be assembled from the exact
  constraint modes of every piece and
  $\tilde{\mathsf S}=\sum_iR_i^\top\Lambda_i[h_i]R_i$ from any fields.  Then
  \begin{enumerate}
    \item $\tilde{\mathsf S}-\mathsf S$ is positive semidefinite;
    \item if $\beta\,|c|^2\le c^\top\mathsf Sc$ for all $c$, then
      $\beta\,|c|^2\le c^\top\tilde{\mathsf S}c$ for all $c$ (with the largest
      such $\beta$: $\lambda_{\min}(\tilde{\mathsf S})\ge\lambda_{\min}(\mathsf S)$);
    \item if $\mathsf S$ is positive definite, so is $\tilde{\mathsf S}$, and
      $\tilde{\mathsf S}c=\chi$ has exactly one solution for every $\chi$.
  \end{enumerate}
\end{theorem}

\begin{proof}
  \leanok
  \uses{cor:T25-ii}
  $\tilde{\mathsf S}-\mathsf S=\sum_iR_i^\top(\Lambda_i[h_i]-\Lambda_i[H_i])R_i$
  is a sum of matrices $R^\top PR$ with $P$ positive semidefinite by
  Corollary~\ref{cor:T25-ii}.  Then
  $c^\top\tilde{\mathsf S}c=c^\top\mathsf Sc+c^\top(\tilde{\mathsf S}-\mathsf S)c
  \ge c^\top\mathsf Sc$, which gives 2 and the definiteness in 3.  A positive
  definite matrix is invertible.
\end{proof}

\section{T25: energy optimality}

\paragraph{In plain words.}
The undivided discrete solution is the admissible field of least total energy.
The host's interface system, formed from the network's fields alone, says
exactly that the rebuilt field has the least energy among the fields that can
be rebuilt from the network's modes.  For a symmetric problem, least energy
over a smaller set is the same as closest to the true minimiser in the energy
size.  So the coupled answer is the best approximation of the undivided
solution that the network's modes can represent: the coupling adds no error of
its own.  Its error is at most the truncation error of the classical
superelement plus the energy of the network's field errors.

\begin{definition}[A quadratic energy]\label{def:quad}
  \lean{Atlas.QuadEnergy, Atlas.QuadEnergy.energy, Atlas.QuadEnergy.norm}
  \leanok
  On a real vector space $X$, let $m$ be a symmetric bilinear form with
  $m(w,w)\ge0$ and $\ell$ a linear functional.  The energy is
  $\mathcal E(w)=\tfrac12m(w,w)-\ell(w)$ and the energy size is
  $\norm{w}_m=\sqrt{m(w,w)}$, a seminorm.
\end{definition}

\begin{lemma}[Minimisers are the stationary points]\label{lem:stationary}
  \lean{Atlas.QuadEnergy.stationary_of_isMin, Atlas.QuadEnergy.isMin_of_stationary,
    Atlas.QuadEnergy.energy_eq_add}
  \leanok
  \uses{def:quad}
  Let $U_0\subseteq X$ be a subspace and $u\in X$.  Then
  $\mathcal E(u)\le\mathcal E(u+v)$ for all $v\in U_0$ if and only if
  $m(u,v)=\ell(v)$ for all $v\in U_0$.  In that case, for every $w$ with
  $w-u\in U_0$,
  \[
    \mathcal E(w)=\mathcal E(u)+\tfrac12\norm{w-u}_m^2 .
  \]
\end{lemma}

\begin{proof}
  \leanok
  $\mathcal E(u+tv)-\mathcal E(u)=t\,\bigl(m(u,v)-\ell(v)\bigr)+\tfrac12t^2m(v,v)$.
  If this is non-negative for all real $t$, the coefficient of $t$ is zero.
  Conversely, if that coefficient is zero the difference is
  $\tfrac12t^2m(v,v)\ge0$; $t=1$ and $v=w-u$ give the identity.
\end{proof}

\begin{lemma}[The abstract lemma: C\'ea with constant one]\label{lem:cea}
  \lean{Atlas.QuadEnergy.closest_iff_isMin}
  \leanok
  \uses{lem:stationary}
  Let $u$ have the least energy in $u+U_0$, let $W\subseteq u+U_0$ be any set
  and $\tilde u\in W$.  Then $\tilde u$ has the least energy in $W$ if and only
  if $\norm{u-\tilde u}_m\le\norm{u-w}_m$ for every $w\in W$.
\end{lemma}

\begin{proof}
  \leanok
  \uses{lem:stationary}
  On $W$, $\mathcal E(w)=\mathcal E(u)+\tfrac12\norm{w-u}_m^2$ by
  Lemma~\ref{lem:stationary}, and $t\mapsto\sqrt t$ is increasing.
\end{proof}

\begin{lemma}[The Galerkin system]\label{lem:galerkin}
  \lean{Atlas.QuadEnergy.galerkin_iff_isMin}
  \leanok
  \uses{lem:stationary}
  For $w_0\in X$, finitely many trial directions $t_a\in X$ and coordinates
  $c$, write $w(c)=w_0+\sum_ac_at_a$.  Then $w(c)$ has the least energy among
  all $w(c')$ if and only if $m\bigl(w(c),t_a\bigr)=\ell(t_a)$ for every $a$.
\end{lemma}

\begin{proof}
  \leanok
  \uses{lem:stationary}
  Lemma~\ref{lem:stationary} with $U_0$ the span of the $t_a$ and $u=w(c)$.
\end{proof}

\begin{lemma}[The energy size obeys the triangle inequality]\label{lem:energy-triangle}
  \lean{Atlas.QuadEnergy.norm_add_le}
  \leanok
  \uses{def:quad}
  $\norm{x+y}_m\le\norm{x}_m+\norm{y}_m$.
\end{lemma}

\begin{proof}
  \leanok
  The Cauchy--Schwarz inequality $m(x,y)^2\le m(x,x)\,m(y,y)$ for a positive
  semidefinite symmetric form (Mathlib:
  \texttt{LinearMap.BilinForm.apply\_sq\_le\_of\_symm}).
\end{proof}

\begin{definition}[The tier-0 solve]\label{def:superelement}
  \lean{Atlas.Superelement, Atlas.Superelement.Admissible, Atlas.Superelement.U0,
    Atlas.Superelement.quad, Atlas.Superelement.IsSolution,
    Atlas.Superelement.coord, Atlas.Superelement.field,
    Atlas.Superelement.portMatrix, Atlas.Superelement.portLoad,
    Atlas.Superelement.hostMatrix, Atlas.Superelement.hostRhs,
    Atlas.Superelement.fieldError}
  \leanok
  \uses{def:gram, def:assemble, def:quad}
  Pieces $i=1,\dots,P$, each with an energy $M_i$ on $U\times F$
  (Definition~\ref{def:piece}), a source $f_i$, port patterns $q_{i,k}\in F$, a
  matrix $R_i$ reading its port coordinates off the free global coordinates
  $c$, and fixed coordinates $r_i$.  A \emph{global field} $w$ gives every
  piece a pair $w_i\in U\times F$.  Let $\Gamma_0$ be a subspace of families of
  face values (those on which neighbours agree and which vanish on the
  prescribed boundary) and $g$ a family carrying the prescribed boundary
  values.  The field $w$ is \emph{admissible} if its face values minus $g$ lie
  in $\Gamma_0$; $U_0$ is the space of fields with face values in $\Gamma_0$.
  The global energy is
  \[
    \mathcal E(w)=\sum_{i=1}^P\Bigl(\tfrac12M_i(w_i,w_i)-f_i(w_{I,i})\Bigr),
    \qquad\norm{w}_M^2=\sum_{i=1}^PM_i(w_i,w_i),
  \]
  and the \emph{undivided solution} $u$ is an admissible field of least energy
  among admissible fields.  From fields $h_{i,k}$ and $\hat u_{p,i}$ the host
  forms
  \begin{gather*}
    \tilde\Lambda_i=\Lambda_i[h_i],\qquad
    \tilde b_{i,k}=f_i(h_{i,k})-M_i\bigl((h_{i,k},q_{i,k}),(\hat u_{p,i},0)\bigr),\\
    \tilde{\mathsf S}=\sum_iR_i^\top\tilde\Lambda_iR_i,\qquad
    \tilde\chi=\sum_iR_i^\top\bigl(\tilde b_i-\tilde\Lambda_ir_i\bigr),
  \end{gather*}
  and for coordinates $c$ rebuilds, with $x_i=r_i+R_ic$,
  \[
    w(c)_i=\Bigl(\hat u_{p,i}+\sum_kx_{i,k}h_{i,k},\ \ \sum_kx_{i,k}q_{i,k}\Bigr).
  \]
\end{definition}

\begin{definition}[The three assumptions]\label{def:assumptions}
  \lean{Atlas.Superelement.Conforming, Atlas.Superelement.BoundaryDataInPortSpan,
    Atlas.Superelement.EnergyNormIsNorm}
  \leanok
  \uses{def:superelement}
  \begin{enumerate}
    \item \emph{Conforming shared sides:} for every $c$, the family
      $\bigl(\sum_k(R_ic)_kq_{i,k}\bigr)_i$ lies in $\Gamma_0$.
    \item \emph{Boundary data in the port span:} the family
      $\bigl(\sum_kr_{i,k}q_{i,k}\bigr)_i-g$ lies in $\Gamma_0$.
    \item \emph{The energy size is a norm on $U_0$:} $w\in U_0$ and
      $\norm{w}_M=0$ imply $w=0$.
  \end{enumerate}
\end{definition}

\begin{theorem}[The host's system is the Galerkin system]\label{thm:T25-galerkin}
  \lean{Atlas.Superelement.hostSystem_iff_isMin}
  \leanok
  \uses{def:superelement, lem:galerkin}
  For any fields $h,\hat u_p$ and any $c$:
  $\tilde{\mathsf S}c=\tilde\chi$ if and only if
  $\mathcal E(w(c))\le\mathcal E(w(c'))$ for every $c'$.
\end{theorem}

\begin{proof}
  \leanok
  \uses{lem:galerkin}
  The trial directions are
  $t_a=\bigl(\sum_k(R_i)_{ka}(h_{i,k},q_{i,k})\bigr)_i$.  Expanding
  $m(w(c),t_a)=\ell(t_a)$ by bilinearity gives row $a$ of
  $\tilde{\mathsf S}c=\tilde\chi$; then Lemma~\ref{lem:galerkin}.
\end{proof}

\begin{theorem}[T25, energy optimality]\label{thm:T25-optimal}
  \lean{Atlas.Superelement.field_admissible, Atlas.Superelement.energy_optimal}
  \leanok
  \uses{thm:T25-galerkin, lem:cea, def:assumptions}
  Assume 1 and 2 of Definition~\ref{def:assumptions}.  Then every rebuilt
  field $w(c)$ is admissible.  If $u$ is the undivided solution and
  $\tilde{\mathsf S}c=\tilde\chi$, then
  \[
    \norm{u-w(c)}_M\le\norm{u-w(c')}_M\qquad\text{for every }c' .
  \]
\end{theorem}

\begin{proof}
  \leanok
  \uses{thm:T25-galerkin, lem:cea}
  The face values of $w(c)$ minus $g$ are the sum of the two families of
  assumptions 1 and 2.  By Theorem~\ref{thm:T25-galerkin}, $w(c)$ has the least
  energy among the rebuilt fields; Lemma~\ref{lem:cea} with $W$ the set of
  rebuilt fields.
\end{proof}

\begin{theorem}[T25, the error bound]\label{thm:T25-bound}
  \lean{Atlas.Superelement.energy_error_bound}
  \leanok
  \uses{thm:T25-optimal, lem:energy-triangle}
  Under the hypotheses of Theorem~\ref{thm:T25-optimal}, for any reference
  fields $H_{i,k}$, $u_{p,i}$ and any coordinates $c^\star$, with
  $x_i^\star=r_i+R_ic^\star$ and $w^{\mathrm{ref}}$ the field rebuilt from the
  reference fields at $c^\star$,
  \[
    \norm{u-w(c)}_M\ \le\ \norm{u-w^{\mathrm{ref}}}_M
      +\Bigl(\sum_{i=1}^Pa_i(d_i,d_i)\Bigr)^{1/2},
    \qquad d_i=(\hat u_{p,i}-u_{p,i})+\sum_kx^\star_{i,k}\,(h_{i,k}-H_{i,k}).
  \]
  With the exact modes and the classical superelement's coordinates,
  $w^{\mathrm{ref}}=u_m$ and this is the document's bound: truncation error
  plus field errors.
\end{theorem}

\begin{proof}
  \leanok
  \uses{thm:T25-optimal, lem:energy-triangle}
  Let $w^\star$ be the field rebuilt from the network's fields at $c^\star$.
  By Theorem~\ref{thm:T25-optimal} and the triangle inequality,
  $\norm{u-w(c)}_M\le\norm{u-w^\star}_M
  \le\norm{u-w^{\mathrm{ref}}}_M+\norm{w^{\mathrm{ref}}-w^\star}_M$.  The field
  $w^\star-w^{\mathrm{ref}}$ has cell values $d_i$ and zero face values on piece
  $i$, and $M_i\bigl((d,0),(d,0)\bigr)=a_i(d,d)$.
\end{proof}

\begin{theorem}[Solvable for every output; one solution]\label{thm:T25-solvable}
  \lean{Atlas.Superelement.solution_unique, Atlas.Superelement.hostMatrix_posDef,
    Atlas.Superelement.hostSystem_existsUnique}
  \leanok
  \uses{def:assumptions, lem:stationary}
  Assume 3 of Definition~\ref{def:assumptions}.  Then the undivided solution
  is unique.  Assume also 1, that each piece's patterns $q_{i,k}$ are linearly
  independent, and that $R_ic=0$ for every $i$ only for $c=0$.  Then
  $\tilde{\mathsf S}$ is positive definite for every choice of the fields $h$,
  and $\tilde{\mathsf S}c=\tilde\chi$ has exactly one solution.
\end{theorem}

\begin{proof}
  \leanok
  \uses{lem:stationary}
  Two solutions $u,u'$ have equal energy, so $\norm{u'-u}_M=0$ by
  Lemma~\ref{lem:stationary}, and $u'-u\in U_0$.  For the matrix,
  $c^\top\tilde{\mathsf S}c=\norm{t(c)}_M^2$ with
  $t(c)_i=\sum_k(R_ic)_k(h_{i,k},q_{i,k})$, which lies in $U_0$ by assumption
  1.  If it vanishes then $t(c)=0$; its face values $\sum_k(R_ic)_kq_{i,k}$
  vanish, so $R_ic=0$ for every $i$, so $c=0$.
\end{proof}

\section{T24: perturbation of the interface solve}

\paragraph{In plain words.}
The exact interface system is $\mathsf Sc=\chi$ and the host solves a learned
one, $\tilde{\mathsf S}\tilde c=\tilde\chi$.  The two answers differ by at most
the error of the matrix times the size of the exact answer, plus the error of
the load, all divided by the smallest stretch $\beta$ of the learned matrix.
For the Gram construction that constant need not be measured on the learned
matrix: by T25 (iii) the exact problem's own constant serves.

\begin{proposition}[T24]\label{prop:T24}
  \lean{Atlas.interface_perturbation}
  \leanok
  Let $\mathsf S,\tilde{\mathsf S}$ be bounded linear maps between real normed
  spaces with $\beta\norm{x}\le\norm{\tilde{\mathsf S}x}$ for all $x$ and some
  $\beta>0$.  If $\mathsf Sc=\chi$ and $\tilde{\mathsf S}\tilde c=\tilde\chi$,
  \[
    \norm{\tilde c-c}\le
      \frac{\norm{\mathsf S-\tilde{\mathsf S}}\,\norm{c}+\norm{\chi-\tilde\chi}}{\beta}.
  \]
\end{proposition}

\begin{proof}
  \leanok
  $\tilde{\mathsf S}(\tilde c-c)=(\tilde\chi-\chi)+(\mathsf S-\tilde{\mathsf S})c$;
  take norms and use the lower bound.
\end{proof}

\begin{corollary}[T24 with the exact problem's constant]\label{cor:T24-dominates}
  \lean{Atlas.interface_perturbation_of_dominates}
  \leanok
  \uses{prop:T24}
  On a real inner-product space, if
  $\beta\norm{x}^2\le\langle\mathsf Sx,x\rangle\le\langle\tilde{\mathsf S}x,x\rangle$
  for all $x$ and some $\beta>0$, the bound of Proposition~\ref{prop:T24} holds
  with this $\beta$.
\end{corollary}

\begin{proof}
  \leanok
  \uses{prop:T24}
  $\beta\norm{x}^2\le\langle\tilde{\mathsf S}x,x\rangle
  \le\norm{\tilde{\mathsf S}x}\,\norm{x}$ by Cauchy--Schwarz.
\end{proof}

\section{T26: the label-free objective}

\paragraph{In plain words.}
The supervised loss compares the network's fields with exact ones, which a
classical solver must first compute.  The label-free objective needs only the
network's own output and the piece's energy.  The two differ by a number fixed
by the piece and its source, so between any two outputs of the network they
change by the same amount: the same gradients and the same minimisers, with no
solved examples.

\begin{definition}[The two objectives]\label{def:objectives}
  \lean{Atlas.Piece.labelFree, Atlas.Piece.supervised}
  \leanok
  \uses{def:gram}
  For returned fields $h_k$ and $\hat u_p$, a source $f$, exact modes $H_k$ and
  the exact particular field $u_p$ (with $a(v,u_p)=f(v)$ for all $v$):
  \[
    \mathcal L_{\mathrm{en}}=\operatorname{tr}\Lambda[h]+a(\hat u_p,\hat u_p)-2f(\hat u_p),
    \qquad
    \mathcal L_{\mathrm{sup}}=\sum_ka(h_k-H_k,h_k-H_k)+a(\hat u_p-u_p,\hat u_p-u_p).
  \]
\end{definition}

\begin{theorem}[T26]\label{thm:T26}
  \lean{Atlas.Piece.labelFree_eq_supervised_add_const,
    Atlas.Piece.labelFree_sub_eq_supervised_sub}
  \leanok
  \uses{def:objectives, cor:T25-ii}
  $\mathcal L_{\mathrm{en}}=\mathcal L_{\mathrm{sup}}
  +\bigl(\operatorname{tr}\Lambda[H]-a(u_p,u_p)\bigr)$.  The bracket does not
  involve the network's output, so for any two outputs
  $\mathcal L_{\mathrm{en}}-\mathcal L_{\mathrm{en}}'
  =\mathcal L_{\mathrm{sup}}-\mathcal L_{\mathrm{sup}}'$.
\end{theorem}

\begin{proof}
  \leanok
  \uses{cor:T25-ii}
  The trace identity of Corollary~\ref{cor:T25-ii}, and
  $a(\hat u_p-u_p,\hat u_p-u_p)=a(\hat u_p,\hat u_p)-2f(\hat u_p)+a(u_p,u_p)$,
  using $a(\hat u_p,u_p)=f(\hat u_p)$.
\end{proof}

\chapter{Tier 2: black-box experts and the contraction contract}

At tier 2 an expert exposes only its boundary response.  The host exchanges
wave variables between neighbouring sides and iterates.  Nothing is automatic
for a learned black box, so the guarantee is conditional: it holds when a
contraction factor has been \emph{certified} for the experts.  This chapter
states what follows from such a certificate (T8, T9a, T9b), and two examples
that show which hypothesis the certificate must be about.

\section{T8: wave variables and the Cayley transform}

\paragraph{In plain words.}
On a side, the effort $e$ (a temperature) and the flow $f$ (a flux into the
piece) are combined into an incoming wave $e+Zf$ and an outgoing wave $e-Zf$,
for an impedance $Z>0$.  The squared size of the outgoing wave is that of the
incoming one minus four times $Z$ times the power the piece absorbs.  So a
piece that never generates power does not amplify waves, and a piece that
absorbs a definite share of what it is given shrinks them by a definite factor.

\begin{lemma}[T8, the wave identity]\label{lem:T8-identity}
  \lean{Atlas.wave_identity}
  \leanok
  In a real inner-product space, for all $e,f$ and all real $Z$,
  \[
    \norm{e-Zf}^2=\norm{e+Zf}^2-4Z\,\langle f,e\rangle .
  \]
\end{lemma}

\begin{proof}
  \leanok
  Expand both squares.
\end{proof}

\begin{theorem}[T8, the bound for any effort and flow]\label{thm:T8}
  \lean{Atlas.wave_bound, Atlas.cayley_bound}
  \leanok
  \uses{lem:T8-identity}
  Let $Z>0$, $c\ge0$, and let $e,f$ satisfy
  $\langle f,e\rangle\ge c\norm{e}^2$ and $\norm{f}\le M\norm{e}$.  Then
  \[
    \norm{e-Zf}^2\le\Bigl(1-\frac{4Zc}{(1+ZM)^2}\Bigr)\norm{e+Zf}^2 .
  \]
  In particular, for a bounded linear $\Lambda$ with
  $\langle\Lambda e,e\rangle\ge c\norm{e}^2$ for all $e$ and
  $\norm{\Lambda}\le M$, this holds with $f=\Lambda e$ for every $e$.
\end{theorem}

\begin{proof}
  \leanok
  \uses{lem:T8-identity}
  With $y=e+Zf$: $\norm{e-Zf}^2=\norm{y}^2-4Z\langle f,e\rangle
  \le\norm{y}^2-4Zc\norm{e}^2$, and $\norm{y}\le(1+ZM)\norm{e}$.
\end{proof}

\begin{theorem}[T8, the scattering map]\label{thm:T8-cayley}
  \lean{Atlas.cayley_transform, Atlas.cayley_nonexpansive}
  \leanok
  \uses{thm:T8}
  Let $E$ be a finite-dimensional real inner-product space, not the zero
  space, and $\Lambda$ as in Theorem~\ref{thm:T8}.  Then $I+Z\Lambda$ is
  invertible and $\mathcal S=(I-Z\Lambda)(I+Z\Lambda)^{-1}$ satisfies
  \[
    \norm{\mathcal S}^2\le1-\frac{4Zc}{(1+ZM)^2}.
  \]
  If only $\langle\Lambda e,e\rangle\ge0$ for all $e$, then $I+Z\Lambda$ is
  invertible and $\norm{\mathcal S}\le1$ (on any finite-dimensional space).
\end{theorem}

\begin{proof}
  \leanok
  \uses{thm:T8}
  $\langle(I+Z\Lambda)e,e\rangle\ge\norm{e}^2$, so $I+Z\Lambda$ is injective,
  hence invertible in finite dimensions.  Every $y$ is $(I+Z\Lambda)e$ for one
  $e$, and $\mathcal Sy=(I-Z\Lambda)e$; Theorem~\ref{thm:T8}.
\end{proof}

\begin{corollary}[T8 at tier 0]\label{cor:T8-gram}
  \lean{Atlas.cayley_matrix_nonexpansive, Atlas.Piece.gram_cayley_nonexpansive}
  \leanok
  \uses{thm:T8, thm:T25-i}
  For a positive semidefinite real matrix $L$ and $Z>0$, $1+ZL$ is invertible
  and $\mathsf S=(1-ZL)(1+ZL)^{-1}$ satisfies $|\mathsf Sa|^2\le|a|^2$ for
  every $a$.  In particular this holds for the Gram port matrix
  $\tilde\Lambda$ of any network output: a tier-0 expert also serves at
  tier 2, with no certificate.
\end{corollary}

\begin{proof}
  \leanok
  \uses{thm:T8, thm:T25-i}
  The case $c=0$, with Theorem~\ref{thm:T25-i}.
\end{proof}

\section{T9a: the contraction contract, one level}

\paragraph{In plain words.}
One sweep applies every expert's scattering map to its incoming wave, hands
each outgoing wave to the neighbouring side, and adds the sources.  If every
expert's map shrinks distances by a certified factor $\rho<1$, and the
hand-over does not lengthen anything, the whole sweep shrinks distances by
$\rho$.  Then T1 applies: one answer, geometric convergence from any start,
and a distance to the classical answer set by how far each learned map is from
the classical one \emph{at the classical answer}.

\begin{definition}[The wave sweep]\label{def:wave-sweep}
  \lean{Atlas.waveSweep}
  \leanok
  Pieces $i$ in a finite set, with wave spaces $W_i$ (normed).  The global wave
  vector $a=(a_i)_i$ is measured in the $\ell^2$ product,
  $d(a,a')^2=\sum_id(a_i,a'_i)^2$.  For maps $S_i:W_i\to W_i$, an exchange
  $\Pi$ on the product and a vector $g$,
  \[
    T(a)=\Pi\bigl((S_i(a_i))_i\bigr)+g .
  \]
\end{definition}

\begin{theorem}[T9a, the sweep contracts]\label{thm:T9a-lip}
  \lean{Atlas.waveSweep_lipschitz}
  \leanok
  \uses{def:wave-sweep}
  If $d(S_ix,S_iy)\le\rho\,d(x,y)$ for every $i$, $x$, $y$, with $\rho\ge0$,
  and $d(\Pi x,\Pi y)\le d(x,y)$, then $d(Ta,Ta')\le\rho\,d(a,a')$.
\end{theorem}

\begin{proof}
  \leanok
  $d(Ta,Ta')^2\le\sum_id(S_ia_i,S_ia'_i)^2\le\rho^2\sum_id(a_i,a'_i)^2$.
\end{proof}

\begin{theorem}[T9a, the contraction contract]\label{thm:T9a}
  \lean{Atlas.contraction_contract}
  \leanok
  \uses{thm:T9a-lip, thm:T1}
  Let the $W_i$ be complete, $S_i$ the learned maps and $S^c_i$ the classical
  ones, with the same $\Pi$ and $g$, and let $a_c$ be a fixed point of the
  classical sweep.  If the hypotheses of Theorem~\ref{thm:T9a-lip} hold with
  $\rho<1$, and $d\bigl(S_i(a_{c,i}),S^c_i(a_{c,i})\bigr)\le\delta_i$ for every
  $i$, then the learned sweep has a unique fixed point $a_l$,
  $d(T^ka_0,a_l)\le\rho^kd(a_0,a_l)$ for every $a_0$ and $k$, and
  \[
    d(a_l,a_c)\le\frac{\bigl(\sum_i\delta_i^2\bigr)^{1/2}}{1-\rho}.
  \]
\end{theorem}

\begin{proof}
  \leanok
  \uses{thm:T9a-lip, thm:T1}
  Theorem~\ref{thm:T1} with the contraction of Theorem~\ref{thm:T9a-lip}; the
  defect at $a_c$ is at most $\bigl(\sum_i\delta_i^2\bigr)^{1/2}$ because
  $\Pi$ does not lengthen distances.
\end{proof}

\section{T9b: two levels}

\paragraph{In plain words.}
With a coarse space, a round is a sweep followed by a coarse solve: a
correction inside the coarse space that leaves no coarse residual.  The first
version of the design asked that each learned map contract on the
\emph{fine} space only, because the coarse space is solved exactly.  That is
not enough (Example~\ref{ex:T9b-diverges}).  What has to be certified is a
contraction factor of the \emph{composite} round.  Given that, the conclusion
is T1's, and the answer is a fixed point of the plain sweep: the coarse solve
changes how the answer is reached, not what it is.

\begin{lemma}[T9b, consistency of the coarse solve]\label{lem:T9b-consistency}
  \lean{Atlas.twoLevel_fixedPt_isFixedPt, Atlas.fixedPt_isFixedPt_twoLevel}
  \leanok
  Let $T$ be the sweep, $C$ the coarse solve and $P_0$ a linear map.  If
  $P_0(Cb-b)=Cb-b$ and $P_0\bigl(Cb-T(Cb)\bigr)=0$ for every $b$, then
  $C(Ta)=a$ implies $Ta=a$.  Conversely, if $Cb=b$ whenever
  $P_0(b-Tb)=0$, then $Ta=a$ implies $C(Ta)=a$.
\end{lemma}

\begin{proof}
  \leanok
  Let $b=Ta$ and $a=Cb$.  Then $a-b$ is fixed by $P_0$, and
  $P_0(a-Ta)=0$; but $a-Ta=a-b$, so $a-b=0$.  The converse is immediate.
\end{proof}

\begin{theorem}[T9b, the two-level contract]\label{thm:T9b}
  \lean{Atlas.two_level_contract}
  \leanok
  \uses{lem:T9b-consistency, thm:T1}
  On a complete normed space, let $T_l,C_l$ be the learned sweep and coarse
  solve, satisfying the hypotheses of Lemma~\ref{lem:T9b-consistency}, and
  $T_c,C_c$ the classical ones, with $C_c(T_ca_c)=a_c$.  If
  $d\bigl(C_lT_lx,C_lT_ly\bigr)\le\rho\,d(x,y)$ for all $x,y$ with
  $0\le\rho<1$, and $d\bigl(C_lT_la_c,C_cT_ca_c\bigr)\le\delta$, then there is
  a unique $a_l$ with $C_l(T_la_l)=a_l$; it satisfies $T_la_l=a_l$;
  $d\bigl((C_lT_l)^ka_0,a_l\bigr)\le\rho^kd(a_0,a_l)$; and
  $d(a_l,a_c)\le\delta/(1-\rho)$.
\end{theorem}

\begin{proof}
  \leanok
  \uses{lem:T9b-consistency, thm:T1}
  Theorem~\ref{thm:T1} for the composite maps, then
  Lemma~\ref{lem:T9b-consistency}.
\end{proof}

\begin{example}[Contracting the fine space is not enough]\label{ex:T9b-diverges}
  \lean{Atlas.two_level_fine_contraction_diverges}
  \leanok
  On $\R^2$ with coarse space the first axis and fine space the second, let
  \[
    G=\begin{pmatrix}0&\tfrac12\\10&0\end{pmatrix},\qquad
    C(x,y)=\bigl(\tfrac12y,\ y\bigr).
  \]
  Then $G$ contracts the fine space by $\tfrac12$; $C$ changes only the coarse
  component and leaves no coarse residual, $\bigl(Ce-G(Ce)\bigr)_1=0$; the only
  fixed point of $G$ is $0$; and
  $(CG)^k(1,2)=\bigl(5^k,\ 2\cdot5^k\bigr)$ for every $k$.
\end{example}

\begin{proof}
  \leanok
  $CG(x,y)=(5x,10x)$.
\end{proof}

\section{A remark beside T1: accuracy is not a substitute}

\paragraph{In plain words.}
T1 asks that the learned map contract.  It cannot be replaced by ``the
classical map contracts and the learned map is uniformly close to it''.  The
example below satisfies both of those, and its learned iteration alternates
between two points for ever.

\begin{example}[The SNI counterexample]\label{ex:sni}
  \lean{Atlas.sni_counterexample}
  \leanok
  \uses{thm:T1}
  Let $T(u)=\tfrac12u$ and $\tilde T(u)=\tfrac12\bigl(u-\tanh10u\bigr)$ on $\R$.
  Then $T$ is a contraction with factor $\tfrac12$ and fixed point $0$;
  $|\tilde T(u)-T(u)|\le\tfrac12$ for every $u$; and there is an
  $x\in[0.3,\,0.34]$ with $\tilde T(x)=-x$ and $\tilde T(-x)=x$, so that
  $\tilde T^k(x)=(-1)^kx$ for every $k$.
\end{example}

\begin{proof}
  \leanok
  $\tilde T(x)=-x$ means $\tanh10x=3x$.  The function $\tanh10x-3x$ is
  continuous, positive at $0.3$ (since $e^6>19$ gives $\tanh3>0.9$) and
  negative at $0.34$ (since $\tanh<1<1.02$); the intermediate value theorem
  gives $x$.  $\tilde T$ is odd.
\end{proof}

